\input{../tex-inputs/latex-leseansicht-vorspann}

\section[Arthur und Olga Schnitzler an Gustav Schwarzkopf, 12. 5. 1914]{L04487 Arthur und Olga Schnitzler an Gustav Schwarzkopf, 12. 5. 1914}
\nopagebreak\mylabel{L04487v}
\rehead{ }\normalsize\beginnumbering\briefempfaengerindex{Schwarzkopf, Gustav@\textsc{Schwarzkopf, Gustav}!zzzSchnitzler, Olga@\emph{von Olga Schnitzler}!1914-05-122@{12. 5. 1914}|(be}\briefempfaengerindex{Schwarzkopf, Gustav@\textsc{Schwarzkopf, Gustav}!zzzSchnitzler, Arthur@\emph{von Arthur Schnitzler}!1914-05-122@{12. 5. 1914}|(be}
\toendnotes[C]{\smallbreak\pagebreak[2]}
\correspDesc{Versand  durch Arthur Schnitzler, Olga Schnitzler am 12. 5. 1914 in Genua
\newline{}Übermittlung  am 13. 5. 1914 in Genua
\newline{}Erhalt  durch Gustav Schwarzkopf im Zeitraum [14. 5. 1914 – 18. 5. 1914?] in Tiefer Graben 17}\toendnotes[C]{\smallbreak}
\Standort{DLA, A:Schnitzler, HS.1985.1.1897.}
\physDesc{Bildpostkarte, 131 Zeichen
\newline{}Handschrift Arthur Schnitzler: Bleistift, deutsche Kurrent
\newline{}Handschrift Olga Schnitzler: Bleistift, lateinische Kurrent
\newline{}Versand: Stempel: »\nobreak{}\oindex{Genua@\textbf{Genua}|pwk}Genova Ferrovia, 13. V. 14, 18\nobreak{}«.  }\pstart{}{\pb}\textsc{Austria\oindex{Österreich@\textbf{Österreich}|pw}}\pend{}\pstart{}Herrn \textsc{Gustav Schwarzkopf}\pend{}\pstart{}Wien I\oindex{I., Innere Stadt@\textbf{I., Innere Stadt}|pw}\pend{}\pstart{}\textsc{Tiefer Graben 17}\oindex{Wien@\textbf{Wien}!I., Innere Stadt@\textbf{I., Innere Stadt}!Tiefer Graben 17@\textbf{Tiefer Graben 17}|pw}\pend{}{\bigskip}
\pstart
           \noindent{}{\pb}\textcolor{gray}{\textbf{Genova – Piazza De Ferrari – Via XX Settembre e nuova
                     Borsa}}\oindex{Palazzo della nuova borsa@\textbf{Palazzo della nuova borsa}|pw}\pend
           \vspace{1em}
\pstart
           {\pb}\substVorne{}\textsuperscript{2}\substDazwischen{}1\substHinten{}2/5 914\pend
           \vspace{0.5em}
\pstart
           {\pb}Herzliche Grüße!\pend
           
\pstart
           Morgen beginnt die Schiffsreiſe.\pend
           \pstart Ihr \spacefill\mbox{Arthur}\pend{}\selectlanguage{ngerman}\vspace{1em}
\pstart
           \noindent{}\raggedleft{}{[}hs. O. Schnitzler:{]} Herzlichst{\\}\spacefill\mbox{Olga}\pend
           \selectlanguage{ngerman}\endnumbering\briefempfaengerindex{Schwarzkopf, Gustav@\textsc{Schwarzkopf, Gustav}!zzzSchnitzler, Olga@\emph{von Olga Schnitzler}!1914-05-122@{12. 5. 1914}|)be}\briefempfaengerindex{Schwarzkopf, Gustav@\textsc{Schwarzkopf, Gustav}!zzzSchnitzler, Arthur@\emph{von Arthur Schnitzler}!1914-05-122@{12. 5. 1914}|)be}\mylabel{L04487h}
\begin{anhang}
\end{anhang}\newcommand{\dateiname}{L04487}\newcommand{\titel}{Arthur und Olga Schnitzler an Gustav Schwarzkopf, 12. 5. 1914}\newcommand{\editorInnen}{Herausgegeben von Jahnke, SelmaMüller, Martin Anton}\input{../tex-inputs/latex-leseansicht-abspann}
